\documentclass[10pt,a4paper]{amsart}
\usepackage[margin=19mm]{geometry}
\usepackage{amsmath,amssymb,mathtools}
\usepackage[scr=rsfs,cal=euler]{mathalfa}
\usepackage{xfrac}
\usepackage[unicode,hidelinks,bookmarks=false]{hyperref}
\newcommand{\ie}{i.e.\ }
\newcommand{\cf}{cf.\ }
\newcommand{\resp}{respectively\ }
\newcommand{\Subsection}{\subsection}
\providecommand{\nobreakdash}{\nobreak\hbox{-}}
\setlength{\emergencystretch}{2em}
\multlinegap0pt
\usepackage{amssymb}
\newtheorem{proposition}{Proposition}
\title{Mod $p$ Monodromy of Cyclic Covers of the Projective Line}
\author{Stepan Nesterov}
\date{Source: arXiv:2604.27391v1 (CC BY 4.0)}
\begin{document}
\maketitle

\noindent\emph{Source and licence.} Mod $p$ Monodromy of Cyclic Covers of the Projective Line. Version arXiv:2604.27391v1 (CC BY 4.0), distributed by arXiv under the Creative Commons Attribution 4.0 International licence, http://creativecommons.org/licenses/by/4.0/, as recorded in the arXiv licence metadata for this exact version and shown on the abstract page https://arxiv.org/abs/2604.27391v1. Full licence notice: \texttt{licenses/arxiv-2604.27391-CC-BY-4.0.txt}.\par\smallskip

\noindent\textit{Editorial scope.} Counted unit: the article's proposition next\_transvect (source0.tex lines 118--120). The article's complete original proof of it is delivered with it (source0.tex lines 121--175, one \texttt{proof} environment of 55 lines). No mathematics is written, repaired or re-derived: the statement and the proof are the article's own lines, in the article's own notation, and no retained line is substituted or re-flowed.  No further source-local context is needed: every cross-reference of every retained line resolves inside the two delivered spans.  Nothing was omitted from any retained span, so the package carries no omission record.  No retained line contains a citation, so the wrapper carries no bibliography and no bibliography entry.  No retained line contains a figure environment, an image-inclusion command or drawing code, so no image is delivered and the package declares no asset dependency.  Editorial formatting only, in addition: an AMS \texttt{amsart} preamble that restates the article's own declarations and macros used by the retained lines, namely the environments \texttt{proposition} on separate counters, together with the standard packages those lines need.  The retained environments print in this document as Proposition 1 in order of appearance, whereas the article prints the same items with the numbering of its own full document.  The complete native source of the article as distributed in the arXiv e-print for this exact version is bundled unchanged as \texttt{sources/199466/source0.tex}.
	\begin{proposition} \label{next_transvect}
		Assume that $q$ is odd and $n > 2$. Let $\bar{g}_{n+1}(t_0, \ldots, t_{n+1})$ be a specialized reduced Gassner representation over $\mathbf{E}_q$ such that $t_0 \ldots t_n = 1$, $t_0 -1 \in \mathbf{E}_q^{\times}$, and $t_{n+1} - 1 \in \mathbf{E}_q^{\times}$. Then there exists an isotropic vector $v$\footnote{This means not $(0,u)$, $(u,0)$} such that the image of $\bar{g}_{n+1}(t_0, \ldots, t_{n+1})$ contains all unitary transvections along $v$.
	\end{proposition}

	\begin{proof} 
		By assumption, the subspace spanned by $\varepsilon_0, \ldots, \varepsilon_{n-1}$, is the space of the degenerate Gassner representation $\bar{g}_n(t_0, \ldots, t_n)$, hence the hermitian form $h$ has a one-dimensional kernel on this hyperplane. We take $v$ to be the spanning vector of the kernel. %Explain better what is going on
		 Let us take an element $T_v(\phi)$ in the unipotent radical of $U(h|_{\langle \varepsilon_0, \ldots, \varepsilon_{n-1} \rangle })$ , and explicitly calculate what are the possible liftings $\widetilde{T}$ of $T_v(\phi)$ to $U(h)$. \par 
		Pick a vector $\hat{\phi}$ such that $\phi(u) = h(u,\hat{\phi})$ for $u \in \langle \varepsilon_0, \ldots, \varepsilon_{n-1} \rangle$. The image $\widetilde{T}(\varepsilon_n)$ has to satisfy $h(u,\widetilde{T}(\varepsilon_n))=h(\widetilde{T}^{-1}(u),\varepsilon_n) = h(u - \phi(u) v, \varepsilon_n)=h(u, \varepsilon_n)-\phi(u)h(v,\varepsilon_n)=h(u,\varepsilon_n)-h(u, \hat{\phi})h(v,\varepsilon_n)=h(u,\varepsilon_n)-h(u, \overline{h(v,\varepsilon_n)}\hat{\phi})$, which implies that $\widetilde{T}(\varepsilon_n)-\varepsilon_n+\overline{h(v,\varepsilon_n)}\hat{\phi}$ can be any multiple of $v$. For convenience, replace $v$ with $\overline{h(v,\varepsilon_n)}^{-1} v$. \par 
		Let us pick a convenient basis of $\bar{g}_n(t_0, \ldots, t_{n+1})$. We will take the first basis element to be $v$, we will let $e_2, \ldots, e_{n}$ to be an orthonormal basis of $\langle \varepsilon_1, \ldots, \varepsilon_m \rangle$, and we will let the last basis element to be $\varepsilon_n$. Then we can write $\hat{\phi} = \alpha_2 e_2 + \ldots \alpha_{n-1} e_{n-1}$, so that $\widetilde{T}$ has the form 
		\begin{equation}\begin{pmatrix}
			1 & \overline{\alpha_2} & \overline{\alpha_3} & \ldots & \overline{\alpha_n} & \lambda \\
			0 & 1 & 0 & \ldots & & \alpha_2 \\
			0 & 0 & 1 & \ldots & & \alpha_3 \\
			\vdots & & & & & \vdots \\
			0 & 0 & 0 & \ldots & & \alpha_n \\
			0 & 0 & 0 & \ldots & & 1
		\end{pmatrix} \label{extension} \end{equation}
		To find unitary transvections with direction $v$ in the image of $\bar{g}_n(t_0, \ldots, t_{n+1})$, consider two such elements $\widetilde{T}_1$, $\widetilde{T}_2$ with distinct first rows $1, \overline{\alpha_2}, \ldots$ and $1, \overline{\beta_2}, \ldots$. For simplicity, assume that $\sum_i \alpha_i \overline{\alpha_i} = \sum_i \beta_i \overline{\beta_i}=1$. Then a direct computation shows that 
		$$\begin{pmatrix}
			1 & \overline{\alpha_2} & \overline{\alpha_3} & \ldots & \overline{\alpha_n} & \lambda \\
			0 & 1 & 0 & \ldots & & \alpha_2 \\
			0 & 0 & 1 & \ldots & & \alpha_3 \\
			\vdots & & & & & \vdots \\
			0 & 0 & 0 & \ldots & & \alpha_n \\
			0 & 0 & 0 & \ldots & & 1
		\end{pmatrix}^{-1} =  \begin{pmatrix}
			1 & -\overline{\alpha_2} & -\overline{\alpha_3} & \ldots & -\overline{\alpha_n} & 1-\lambda \\
			0 & 1 & 0 & \ldots & & -\alpha_2 \\
			0 & 0 & 1 & \ldots & & -\alpha_3 \\
			\vdots & & & & & \vdots \\
			0 & 0 & 0 & \ldots & & -\alpha_n \\
			0 & 0 & 0 & \ldots & & 1
		\end{pmatrix}$$
		and 
		$$\Big[ \begin{pmatrix}
			1 & \overline{\alpha_2} & \overline{\alpha_3} & \ldots & \overline{\alpha_n} & \lambda \\
			0 & 1 & 0 & \ldots & & \alpha_2 \\
			0 & 0 & 1 & \ldots & & \alpha_3 \\
			\vdots & & & & & \vdots \\
			0 & 0 & 0 & \ldots & & \alpha_n \\
			0 & 0 & 0 & \ldots & & 1
		\end{pmatrix} 
		\begin{pmatrix}
			1 & \overline{\beta_2} & \overline{\beta_3} & \ldots & \overline{\beta_n} & \mu \\
			0 & 1 & 0 & \ldots & & \beta_2 \\
			0 & 0 & 1 & \ldots & & \beta_3 \\
			\vdots & & & & & \vdots \\
			0 & 0 & 0 & \ldots & & \beta_n \\
			0 & 0 & 0 & \ldots & & 1
		\end{pmatrix} \Big] = \begin{pmatrix}
			1 & 0 & 0 & \ldots & 0 & \sum_i \overline{\alpha_i}\beta_i - \alpha_i \overline{\beta_i} \\
			0 & 1 & 0 & \ldots & & 0 \\
			0 & 0 & 1 & \ldots & & 0 \\
			\vdots & & & & & \vdots \\
			0 & 0 & 0 & \ldots & & 0 \\
			0 & 0 & 0 & \ldots & & 1
		\end{pmatrix}$$
		I claim that for $n > 2$, we may pick $\widetilde{T}_1$ and $\widetilde{T}_2$ in such a way that $ \sum_i \overline{\alpha_i}\beta_i - \alpha_i \overline{\beta_i} $ is any imaginary scalar. Indeed, let $\xi$ be an imaginary scalar. Because the norm map is surjective, there is an $\eta$ such that $\frac{\xi \overline{\xi}}{4} + \eta \overline{\eta}=1$. Then we take $(\alpha_2, \alpha_3, \ldots, \alpha_n) = (\frac{\xi}{2}, \eta, 0,\ldots, 0)$, $(\beta_2, \beta_3, \ldots, \beta_n) = (1,0, \ldots, 0)$.
	\end{proof}

\end{document}
